\documentclass[a4paper, 14pt]{extarticle}

% Поля
%--------------------------------------
\usepackage{geometry}
\geometry{a4paper,tmargin=2cm,bmargin=2cm,lmargin=3cm,rmargin=1cm}
%--------------------------------------


%Russian-specific packages
%--------------------------------------
\usepackage[T2A]{fontenc}
\usepackage[utf8]{inputenc} 
\usepackage[english, main=russian]{babel}
%--------------------------------------

\usepackage{textcomp}

% Красная строка
%--------------------------------------
\usepackage{indentfirst}               
%--------------------------------------             


%Graphics
%--------------------------------------
\usepackage{graphicx}
\graphicspath{ {./images/} }
\usepackage{wrapfig}
\usepackage{listings}
%--------------------------------------

% Полуторный интервал
%--------------------------------------
\linespread{1.3}                    
%--------------------------------------

%Выравнивание и переносы
%--------------------------------------
% Избавляемся от переполнений
\sloppy
% Запрещаем разрыв страницы после первой строки абзаца
\clubpenalty=10000
% Запрещаем разрыв страницы после последней строки абзаца
\widowpenalty=10000
%--------------------------------------

%Списки
\usepackage{enumitem}

%Подписи
\usepackage{caption} 

%Гиперссылки
\usepackage{hyperref}

\usepackage{minted}

\hypersetup {
	unicode=true
}

%Рисунки
%--------------------------------------
\DeclareCaptionLabelSeparator*{emdash}{~--- }
\captionsetup[figure]{labelsep=emdash,font=onehalfspacing,position=bottom}
%--------------------------------------

\usepackage{tempora}
\usepackage{amsmath}
\usepackage{color}
\usepackage{listings}
\lstset{
  belowcaptionskip=1\baselineskip,
  breaklines=true,
  frame=L,
  xleftmargin=\parindent,
  language=Python,
  showstringspaces=false,
  basicstyle=\footnotesize\ttfamily,
  keywordstyle=\bfseries\color{blue},
  commentstyle=\itshape\color{purple},
  identifierstyle=\color{black},
  stringstyle=\color{red},
  literate={а}{{\selectfont\char224}}1
         {б}{{\selectfont\char225}}1
         {в}{{\selectfont\char226}}1
         {г}{{\selectfont\char227}}1
         {д}{{\selectfont\char228}}1
         {е}{{\selectfont\char229}}1
         {ё}{{\"e}}1
         {ж}{{\selectfont\char230}}1
         {з}{{\selectfont\char231}}1
         {и}{{\selectfont\char232}}1
         {й}{{\selectfont\char233}}1
         {к}{{\selectfont\char234}}1
         {л}{{\selectfont\char235}}1
         {м}{{\selectfont\char236}}1
         {н}{{\selectfont\char237}}1
         {о}{{\selectfont\char238}}1
         {п}{{\selectfont\char239}}1
         {р}{{\selectfont\char240}}1
         {с}{{\selectfont\char241}}1
         {т}{{\selectfont\char242}}1
         {у}{{\selectfont\char243}}1
         {ф}{{\selectfont\char244}}1
         {х}{{\selectfont\char245}}1
         {ц}{{\selectfont\char246}}1
         {ч}{{\selectfont\char247}}1
         {ш}{{\selectfont\char248}}1
         {щ}{{\selectfont\char249}}1
         {ъ}{{\selectfont\char250}}1
         {ы}{{\selectfont\char251}}1
         {ь}{{\selectfont\char252}}1
         {э}{{\selectfont\char253}}1
         {ю}{{\selectfont\char254}}1
         {я}{{\selectfont\char255}}1
         {А}{{\selectfont\char192}}1
         {Б}{{\selectfont\char193}}1
         {В}{{\selectfont\char194}}1
         {Г}{{\selectfont\char195}}1
         {Д}{{\selectfont\char196}}1
         {Е}{{\selectfont\char197}}1
         {Ё}{{\"E}}1
         {Ж}{{\selectfont\char198}}1
         {З}{{\selectfont\char199}}1
         {И}{{\selectfont\char200}}1
         {Й}{{\selectfont\char201}}1
         {К}{{\selectfont\char202}}1
         {Л}{{\selectfont\char203}}1
         {М}{{\selectfont\char204}}1
         {Н}{{\selectfont\char205}}1
         {О}{{\selectfont\char206}}1
         {П}{{\selectfont\char207}}1
         {Р}{{\selectfont\char208}}1
         {С}{{\selectfont\char209}}1
         {Т}{{\selectfont\char210}}1
         {У}{{\selectfont\char211}}1
         {Ф}{{\selectfont\char212}}1
         {Х}{{\selectfont\char213}}1
         {Ц}{{\selectfont\char214}}1
         {Ч}{{\selectfont\char215}}1
         {Ш}{{\selectfont\char216}}1
         {Щ}{{\selectfont\char217}}1
         {Ъ}{{\selectfont\char218}}1
         {Ы}{{\selectfont\char219}}1
         {Ь}{{\selectfont\char220}}1
         {Э}{{\selectfont\char221}}1
         {Ю}{{\selectfont\char222}}1
         {Я}{{\selectfont\char223}}1
}

%--------------------------------------
%			НАЧАЛО ДОКУМЕНТА
%--------------------------------------

\begin{document}

%--------------------------------------
%			ТИТУЛЬНЫЙ ЛИСТ
%--------------------------------------
\begin{titlepage}
\thispagestyle{empty}
\newpage


%Шапка титульного листа
%--------------------------------------
\vspace*{-60pt}
\hspace{-65pt}
\begin{minipage}{0.3\textwidth}
\hspace*{-20pt}\centering
\includegraphics[width=\textwidth]{emblem}
\end{minipage}
\begin{minipage}{0.67\textwidth}\small \textbf{
\vspace*{-0.7ex}
\hspace*{-6pt}\centerline{Министерство науки и высшего образования Российской Федерации}
\vspace*{-0.7ex}
\centerline{Федеральное государственное бюджетное образовательное учреждение }
\vspace*{-0.7ex}
\centerline{высшего образования}
\vspace*{-0.7ex}
\centerline{<<Московский государственный технический университет}
\vspace*{-0.7ex}
\centerline{имени Н.Э. Баумана}
\vspace*{-0.7ex}
\centerline{(национальный исследовательский университет)>>}
\vspace*{-0.7ex}
\centerline{(МГТУ им. Н.Э. Баумана)}}
\end{minipage}
%--------------------------------------

%Полосы
%--------------------------------------
\vspace{-25pt}
\hspace{-35pt}\rule{\textwidth}{2.3pt}

\vspace*{-20.3pt}
\hspace{-35pt}\rule{\textwidth}{0.4pt}
%--------------------------------------

\vspace{1.5ex}
\hspace{-35pt} \noindent \small ФАКУЛЬТЕТ\hspace{80pt} <<Информатика и системы управления>>

\vspace*{-16pt}
\hspace{47pt}\rule{0.83\textwidth}{0.4pt}

\vspace{0.5ex}
\hspace{-35pt} \noindent \small КАФЕДРА\hspace{50pt} <<Теоретическая информатика и компьютерные технологии>>

\vspace*{-16pt}
\hspace{30pt}\rule{0.866\textwidth}{0.4pt}
  
\vspace{11em}

\begin{center}
\Large {\bf Лабораторная работа №1} \\ 
\large {\bf по курсу <<Базы данных>>} \\ 
\large <<Моделирование данных с использованием модели сущность-связь>>
\end{center}\normalsize

\vspace{8em}


\begin{flushright}
  {Студент группы ИУ9-41Б Росоха М. С.\hspace*{15pt} \\
  \vspace{2ex}
  Преподаватель Вишняков И. Э.\hspace*{15pt}}
\end{flushright}

\bigskip

\vfill
 

\begin{center}
\textsl{Москва 2026}
\end{center}
\end{titlepage}
%--------------------------------------
%		КОНЕЦ ТИТУЛЬНОГО ЛИСТА
%--------------------------------------

\renewcommand{\ttdefault}{pcr}

\setlength{\tabcolsep}{3pt}
\newpage
\setcounter{page}{2}
\section{Задачи}\label{Sect::task}
\indent • Выбрать простейшую предметную область, соответствующую 4-5 сущностям. \newline
\indent • Сформировать требования к предметной области. \newline
\indent • Создать модель «сущность-связь» для предметной области с обоснованием выбора кардинальных чисел связей.
\section{Практическая реализация}
В качестве предметной области была выбрана сфера бухгалтерского учёта в разделе кадров в некоторой фирме. Данная схема (рис. 1) удовлетворяет следующим функциональным требованиям:
\begin{itemize}
    \item Учёт сотрудников, когда-либо работавших в данной компании.
    \item Формирование финансовых отчётов по указанному периоду и отделу c учётом отпускных, больничных и премий.
    \item Учёт отделов с поддержкой иерархии.
\end{itemize}

\begin{figure}[h]
    \centering
    \includegraphics[width=0.9\textwidth]{Снимок экрана 2026-09-14 115630.png}
    \caption{ER-схема базы данных}
    \label{fig:er}
\end{figure}
\newpage
\noindent В данной ER-схеме используются следующие сущности:
\begin{description}
    \item[1. Сотрудник] хранит основную информацию о человеке, работавшем в фирме.\newline
    Атрибуты сущности:
    \vspace{-0.5em}
    \begin{itemize}
        \item ИНН - ID
        \item СНИЛС
        \item Фамилия, имя, отчество
        \item Дата рождения
        \item Серия и номер паспорта
        \item Телефон
        \item Почта
    \end{itemize}
    
    \item[2. Договор] содержит всю необходимую информацию о трудовом договоре.\newline
    Атрибуты сущности:
    \vspace{-0.5em}
    \begin{itemize}
        \item Номер договора - ID
        \item ИНН
        \item Оклад
        \item Часы работы
        \item Дата подписания
        \item Дата начала работы
        \item Дата увольнения
        \item Причина увольнения
    \end{itemize}
    
    \item[3. Отдел] хранит общие данные о подразделении и отображает их структуру.\newline
    Атрибуты сущности:
    \vspace{-0.5em}
    \begin{itemize}
        \item Название отдела - ID
        \item Бюджет
        \item Численность (штат)
    \end{itemize}
    
    \item[4. Должность] отображает основную информацию о должности в компании.\newline
    Атрибуты сущности:
    \vspace{-0.5em}
    \begin{itemize}
        \item Название должности - ID
        \item Минимальный оклад
        \item Максимальный оклад
        \item Норма часов в неделю
    \end{itemize}
    
    \item[5. Отпуск] показывает необходимую информацию об отпусках сотрудников.\newline
    Атрибуты сущности: 
    \vspace{-0.5em}
    \begin{itemize}
        \item Номер приказа - ID
        \item ИНН - ID
        \item Дата начала
        \item Дата окончания
        \item Сумма отпускных
    \end{itemize}
    
    \item[6. Больничный] фиксирует данные больничного сотрудника. \newline
    Атрибуты сущности:
    \vspace{-0.5em}
    \begin{itemize}
        \item Номер листа - ID
        \item ИНН
        \item Дата начала
        \item Дата окончания
        \item Сумма выплат
    \end{itemize}

    \item[7. Премия] хранит информацию о различных премиях, выданных компанией.\newline
    Атрибуты сущности:
    \vspace{-0.5em}
    \begin{itemize}
        \item Номер приказа - ID
        \item ИНН - ID
        \item Тип премии
        \item Сумма
        \item Дата выдачи
    \end{itemize}
\end{description}
\newpage
\indent Между указанными сущностями используются следующие связи:
\begin{enumerate}
    \item Сотрудник и договор связаны отношением один ко многим (максимальные кардинальные числа - 1 и M, минимальные кардинальные числа - 1 и 1): у сотрудника может быть несколько договоров (несколько должностей в одно время или в течение опредлённого времени), а у договоров только один сотрудник.
    \vspace{-0.25em}
    \item Сотрудник и отпуск связаны отношением один ко многим (максимальные кардинальные числа - M и 1, минимальные кардинальные числа - 0 и 1): у сотрудника может быть ноль или несколько отпусков за время работы, но каждый отпуск принадлежит только одному сотруднику.
    \vspace{-0.25em}
    \item Сотрудник и больничный связаны отношением один ко многим (максимальные кардинальные числа - M и 1, минимальные кардинальные числа - 0 и 1): у сотрудника может быть ноль или несколько больничных, но каждый больничный оформлен на одного сотрудника.
    \vspace{-0.25em}
    \item Сотрудник и премия связаны отношением один ко многим (максимальные кардинальные числа - M и 1, минимальные кардинальные число - 0 и 1): у сотрудника может быть ноль или несколько премий, но каждая премия выдана только одному сотруднику.
    \vspace{-0.25em}
    \item Договор и отдел связаны отношением многие к одному (максимальные кардинальные числа - 1 и М, минимальные кардинальные числа - 1 и 0): каждый договор заключается в одном отделе, но в одном отделе может быть заключено ноль (отдел новый) или несколько договоров.
    \item Договор и должность связаны отношением многие к одному (максимальные кардинальные числа - 1 и M, минимальные кардинальные числа - 1 и 0): каждый договор соответствует одной должности, но на одну должность может быть принято ноль (новая должность) или несколько сотрудников.
    \vspace{-0.25em}
    \item Отдел связан с самим собой рекурсивной связью многие к одному (максимальные кардинальные числа - M и 1, минимальные кардинальные числа - 0 и 0): у каждого отдела может быть один вышестоящий отдел или не быть вовсе, но один вышестоящий отдел может включать несколько подчинённых или не иметь одного.
\end{enumerate}
\end{document}